\subsection{Word2Vec}\label{sec:word2vec}

\subsubsection{From Neural Language Models to Word2Vec}\label{sec:from-neural-language-models-to-word2vec}

At the end of the previous section we identified the main paradox of Bengio's Neural Language Model: \hl{it solved the curse of dimensionality, but it was too computationally expensive to scale.} This is where \textbf{Word2Vec} enters the story.\cite{mikolov2013efficient} In 2013, a team at Google led by \textbf{Tomas Mikolov} asked a simple but revolutionary question: \emph{do we really need a full language model to learn good word embeddings?} Their answer was a \textbf{no}, and they focused directly on learning the word vectors, without the need to learn a full language model.

\highspace
\textcolor{Green3}{\faIcon[regular]{lightbulb} \textbf{The New Idea.}} Mikolov noticed something interesting. When people used Bengio's model, what they really cared about was often not the language model itself. They cared about the learned word vectors. The embeddings were the valuable product. The neural language model was merely a mechanism to obtain them. Therefore, instead of training a large language model to obtain embeddings as a side effect, \hl{Word2Vec proposes to directly train specifically to learn embeddings, and forget about building a complex Language Model.} This shift completely changed the design philosophy. The quote by linguist John Rupert Firth (1957), \emph{``You shall know a word by the company it keeps''}, became the guiding principle of Word2Vec. The idea was to \hl{learn word embeddings by predicting a word based on its context, or vice versa, without the need for a full language model.}

\highspace
For example, imagine a corpus containing sentences such as:
\begin{itemize}
    \item \emph{Pelé has called Neymar an excellent player}
    \item \emph{At the age of 22 years, Neymar had scored 40 goals}
    \item \emph{Neymar was called a true phenomenon}
\end{itemize}
The word \emph{Neymar} appears surrounded by words like ``\emph{player}'', ``\emph{goals}'', ``\emph{phenomenon}'', ``\emph{Pelé}'', ``\emph{football}''. Word2Vec learns the representation of \emph{Neymar} from these neighboring words. The idea is that the context defines the meaning of a word.

\highspace
\textcolor{Green3}{\faIcon{tachometer-alt} \textbf{Efficiency Improvements.}} The biggest contribution of Word2Vec was not necessarily better \textbf{embeddings}. It was simply a \textbf{much faster way to learn} them.
\begin{itemize}
    \item \textbf{Bye Bye Hidden Layer.} Bengio's model contained a hidden layer, which was computationally expensive. The idea of Mikolov was to \hl{remove the hidden layer entirely}, and make the architecture much shallower.
    

    \item \textbf{Shared Projection Layer}. In Bengio's model, each word had its own projection layer. Each context position had its own contribution to the hidden layer. Word2Vec instead focuses on learning a single embedding space used throughout the model. Therefore, the emphasis shifts from \emph{learning a complex predictor} to \textbf{learning high-quality vectors} (the embeddings). The model is now much simpler, and the training is much faster.
    

    \item \textbf{Large Contexts}. Another important difference is that Word2Vec uses words from the past and future as context, instead of just the previous words. This allows the model to learn from a larger context window, which can lead to better embeddings.
\end{itemize}
In general, the Bengio's goal was to get the best possible language model, while Mikolov's goal was to get the \textbf{best possible word vectors}. These are related but not identical objectives.

\highspace
In general, the main idea was to achieve better performance by allowing a simpler (shallower) model to be trained on much larger amounts of data, but without the need to learn a full language model. And this is one of the most important lessons in machine learning: \hl{a simpler model trained on much more data can outperform a sophisticated model trained on less data.} Word2Vec embraced this philosophy completely.